\subsection{How far from the ground is a state of given energy?}
\label{subsec:fidelity-energy}

Equation~\eqref{eq:ground-return-energy} was stated for the partner $r$ of the output ground, in the domain of $H_\Psi$. The same identity holds for every state of finite energy in either sector, and with a gap it bounds the distance of the state from the ground ray. Let $H$ be self-adjoint and bounded below on a Hilbert space, let $E_0$ be the bottom of its spectrum, and let $x$ be a unit vector with $Hx=E_0x$. Write $P_x$ for its projection, $Q_x=I-P_x$, and for a unit vector $\psi$
\begin{equation}
 F_x(\psi)=|\langle x,\psi\rangle|^2,\qquad
 d_{\rm FS}([\psi],[x])=\arccos|\langle x,\psi\rangle|,\qquad
 1-F_x(\psi)=\sin^2d_{\rm FS}([\psi],[x]).
 \label{eq:fidelity-energy-distance}
\end{equation}
The number $F_x$ of Eq.~\eqref{eq:actual-ground-partner} is $F_x(r)$, and $d_{\rm FS}$ is the Fubini--Study distance of Section~\ref{sec:polar-matched-observables}, measured here from the ray of $x$. The excitation energy is defined on the form domain $\operatorname{Dom}(H-E_0I)^{1/2}$, which contains $\operatorname{Dom}H$:
\begin{equation}
 E_{\rm exc}(\psi)=\bigl\|(H-E_0I)^{1/2}\psi\bigr\|^2,
 \qquad E_{\rm exc}(\psi)=\langle\psi,(H-E_0I)\psi\rangle\ \text{ for }\psi\in\operatorname{Dom}H .
 \label{eq:fidelity-energy-form}
\end{equation}
By the spectral theorem it is nonnegative on the whole form domain: no state has energy below $E_0$. This uses only that $E_0$ is the bottom of the spectrum, not that it is an eigenvalue. Uniqueness of the ground ray is a further assumption, as in R2.

\begin{proposition}[Energy and fidelity]
\label{prop:fidelity-energy}
Let $\psi$ be a unit vector in the form domain.
\begin{enumerate}
\item[(a)] The vector $Q_x\psi$ lies in the form domain and $E_{\rm exc}(\psi)=E_{\rm exc}(Q_x\psi)$. If $F_x(\psi)<1$, put $\eta_\psi=Q_x\psi/\sqrt{1-F_x(\psi)}$. Then $\eta_\psi$ is a unit vector orthogonal to $x$, and
\begin{equation}
 E_{\rm exc}(\psi)=\bigl(1-F_x(\psi)\bigr)\,Z(\psi),\qquad Z(\psi)=E_{\rm exc}(\eta_\psi)\ge0 .
 \label{eq:fidelity-energy-factor}
\end{equation}
If $F_x(\psi)=1$, then $E_{\rm exc}(\psi)=0$.
\item[(b)] Suppose that $E_{\rm exc}(\phi)\ge\Delta\,\|Q_x\phi\|^2$ for every $\phi$ in the form domain, with some $\Delta>0$. Then $Z(\psi)\ge\Delta$ and
\begin{equation}
 E_{\rm exc}(\psi)\ge\Delta\,\bigl(1-F_x(\psi)\bigr)=\Delta\,\sin^2d_{\rm FS}([\psi],[x]).
 \label{eq:fidelity-energy-gap}
\end{equation}
In particular $E_{\rm exc}(\psi)=0$ only on the ray of $x$.
\item[(c)] If $H$ is bounded with top of spectrum $E_{\max}$, then $Z(\psi)\le E_{\max}-E_0$ and $E_{\rm exc}(\psi)\le(E_{\max}-E_0)\bigl(1-F_x(\psi)\bigr)$.
\end{enumerate}
\end{proposition}
\begin{proof}
The vector $x$ lies in $\operatorname{Dom}H$ and $\|(H-E_0I)^{1/2}x\|^2=\langle x,(H-E_0I)x\rangle=0$. Hence $Q_x\psi=\psi-\langle x,\psi\rangle x$ is in the form domain and $(H-E_0I)^{1/2}Q_x\psi=(H-E_0I)^{1/2}\psi$. This is the first statement of (a), and $\|Q_x\psi\|^2=1-F_x(\psi)$ gives the rest. For (b) apply the hypothesis to $\psi$ and to $\eta_\psi$. For (c) use $0\le H-E_0I\le(E_{\max}-E_0)I$ on $\eta_\psi$.
\end{proof}

The hypothesis of (b) says that the spectrum of $H$ on $x^\perp$ begins at least $\Delta$ above $E_0$. It makes the ground ray unique. The number $Z(\psi)$ is the mean excitation energy of $\eta_\psi$, the state that an ideal projective ground test leaves when $\psi$ fails it. It depends on the state, and it is constant only under further spectral structure, for instance when $H-E_0I$ is a multiple of the identity on $x^\perp$. For $0<F_x(\psi)<1$ there is a phase $\alpha$ with
\begin{equation}
 \psi=e^{i\alpha}\sqrt{F_x(\psi)}\;x+\sqrt{1-F_x(\psi)}\;\eta_\psi .
 \label{eq:fidelity-energy-decomposition}
\end{equation}
At $F_x(\psi)=0$ the first term is absent and no phase is needed. At $F_x(\psi)=1$ the vector $\eta_\psi$ is not defined, and part (a) does not use it.

Three cases show the range of $Z$. On the two-site path of Section~\ref{sec:ground-correspondence} the complement of the ground is one-dimensional, so $Z$ equals the gap $1$ for every state. The partner $(1,2)/\sqrt5$ has $E_{\rm exc}=1-F_x=1/10$. On the five-site path the gap is $1-\cos(\pi/5)\approx0.19$ and the top of the spectrum is $1+\cos(\pi/5)\approx1.81$. The site state $e_2$ of Section~\ref{sec:determination} has fidelity $1/5$ with the constant ground, excitation energy $1$ and $Z=5/4$. Third, the form domain is larger than the operator domain when $H$ is unbounded, and the proposition applies there. For the number operator $N$ of Section~\ref{sec:four-maps}, with ground $e_0$ and gap one, let $\psi$ have probabilities proportional to $(n+1)^{-5/2}$ on the vectors $e_n$. Its mean energy is finite and its second moment is infinite, so it lies in the form domain and outside $\operatorname{Dom}N$. In terms of the zeta function,
\begin{equation}
 F_{e_0}(\psi)=\frac1{\zeta(5/2)}\approx0.745,\qquad
 E_{\rm exc}(\psi)=\frac{\zeta(3/2)}{\zeta(5/2)}-1\approx0.947,\qquad Z(\psi)\approx3.72 .
 \label{eq:fidelity-energy-oscillator}
\end{equation}

Without the bound of (c) a fidelity close to one does not control the energy. If $H$ is unbounded there are, for large $k$, unit vectors $\phi_k$ orthogonal to $x$ with $E_{\rm exc}(\phi_k)=k^2$. The states $\sqrt{1-1/k}\;x+\phi_k/\sqrt k$ have fidelity $1-1/k$ and excitation energy $k$. In the other direction a gap suffices: by Eq.~\eqref{eq:fidelity-energy-gap}, states whose excitation energy tends to zero converge to the ray of $x$ in the Fubini--Study distance.

Convergence to the ground ray can also be read from the components. A nonzero vector $cx+\eta$ with $\eta\perp x$ and $c\ne0$ has, after normalization,
\begin{equation}
 F_x=\frac{|c|^2}{|c|^2+\|\eta\|^2}=\frac1{1+\|\eta\|^2/|c|^2}.
 \label{eq:fidelity-energy-ratio}
\end{equation}
A sequence of such vectors converges to the ray of $x$ exactly when $\|\eta_n\|/|c_n|\to0$. For determinations $X_n$ acting on the ground $y$ of Section~\ref{sec:determination}, $c_n=\omega(X_n)$ and $\eta_n=X_n^\circ y$. The states $X_ny/\|X_ny\|$ therefore return to the ground ray exactly when $\|X_n^\circ y\|/|\omega(X_n)|\to0$.

The proposition concerns one Hamiltonian and one ground. It applies to the input with $H_\Psi$ and $x$, and to the output with $H_\Phi$ and $y$, separately. The polar map relates the fidelities of the two sides and not their energies. Under the weak bridge, for a unit vector $\psi$ in $S$,
\begin{equation}
 F_y(V\psi)=F_r(\psi),\qquad r=u^*y,
 \label{eq:fidelity-energy-transport}
\end{equation}
with $r$ the partner of Eq.~\eqref{eq:actual-ground-partner}. Under strong compatibility this is $F_x(\psi)$. The graph maps $G_\varphi$ of Section~\ref{sec:neutral-graphs-phase} are isometries as well, so for unit vectors $\psi,\chi$ in $S$ the rays of $\psi$, $V\psi$ and $G_\varphi\psi$ have the same Fubini--Study distances from those of $\chi$, $V\chi$ and $G_\varphi\chi$. The energies of $\psi$ and $V\psi$ are not related in this way unless $V$ intertwines the two Hamiltonians, which R1--R5 do not require. On the unequal-cell path the coarse ground has energy zero and its image under $V$ has energy $1/16$ on the fine path.

The proposition needs a self-adjoint operator bounded below and a ground vector, and nothing else. It can be applied to a supplied sector Hamiltonian, to the return $A^*A$, or to the effective operator $H_P-C^*(H_Q-z)^{-1}C$ of Eq.~\eqref{eq:prepared-schur} at a real $z$ in the resolvent set of $H_Q$, where it is self-adjoint. The hypotheses have to be checked in each case. For $A^*A$ the form domain is $\operatorname{Dom}A$ and the form is $\|A\psi\|^2$. If $x$ spans $\ker A$, then $E_0=0$ and $E_{\rm exc}(\psi)=\|A\psi\|^2$, the hypothesis of (b) is the lower bound of Section~\ref{sec:cohomology-index} that makes the range of $A$ closed, and the image $Ax$ is zero. A reference vector with nonzero image is in general not a ground of $A^*A$. Under strong compatibility it is an eigenvector of $A^*A$ with eigenvalue $a^2$, which need not be the bottom of the spectrum. No cyclicity is used.
